\ifdefined\gkver\else\def\gkver{student}\fi
\documentclass[\gkver]{gaokaozhenti}
\begin{document}

\chapter{2000年天津}
%% paperType: 理综物理部分

\begin{enumerate}
\item
%% number: 3
%% paperName: 2000年天津高考第 3 题 4 分
%% typeId: 1
%% type: 单选题
%% chapter: 曲线运动
%% point: 圆周运动的应用
%% method: 无
%% score: 4
%% degree: 550
%% duplicateId: 0
%% body:
在高速公路的拐弯处，路面造得外高内低，即当车向右拐弯时，司机左侧的路面比右侧的要高一些，路面与水平面间的夹角为 $\theta$。设拐弯路段是半径为 $R$ 的圆弧，要使车速为 $v$ 时车轮与路面之间的横向（即垂直于前进方向）摩擦力等于零，$\theta$ 应等于 \xzanswer{B}

\fourchoices[answer=B]
{$\arcsin\frac{v^{2}}{Rg}$}
{$\arctan\frac{v^{2}}{Rg}$}
{$\frac{1}{2}\arcsin\frac{2v^{2}}{Rg}$}
{$\mathrm{arccot}\frac{v^{2}}{Rg}$}

%% answer: B
%% memo:
\memoanswer{本题原卷及材料中未提供详解，待补充。}

\item
%% number: 9
%% paperName: 2000年天津高考第 9 题 4 分
%% typeId: 1
%% type: 单选题
%% chapter: 电路
%% point: 动态分析
%% method: 无
%% score: 4
%% degree: 550
%% duplicateId: 0
%% body:
如图~\ref{2000天津09} 所示，A、B、C 是三相交流电源的三根相线，O 是中线，电源的相电压为 $220\UV$，$L_{1}$、$L_{2}$、$L_{3}$ 是三个“$220\UV$ $60\UW$”的灯泡，开关 $K_{1}$ 断开，$K_{2}$、$K_{3}$ 闭合，由于某种原因，电源中线在图中 $O'$ 处断了，那么 $L_{2}$ 和 $L_{3}$ 两灯泡将 \xzanswer{C}

\onepicture[num=9, showlabel=false, label=2000天津09, w=6cm]{figs/09.png}

\fourchoices[answer=C]
{立刻熄灭}
{变得比原来亮一些}
{变得比原来暗一些}
{保持亮度不变}

%% answer: C
%% memo:
\memoanswer{本题原卷及材料中未提供详解，待补充。}

\item
%% number: 11
%% paperName: 2000年天津高考第 11 题 5 分
%% typeId: 3
%% type: 填空题
%% chapter: 光学
%% point: 光电效应
%% method: 无
%% score: 5
%% degree: 750
%% duplicateId: 0
%% body:
已知金属铯的逸出功为 $1.9\UeV$，在光电效应实验中，要使铯表面发出的光电子的最大动能为 $1.0\UeV$，入射光的波长应为 \tkanswer{$4.3\times10^{-9}$}$\Um$。

%% answer: 4.3\times10^{-9}
\jdanswer{$4.3\times10^{-9}\Um$}

%% memo:
\memoanswer{本题原卷及材料中未提供详解，待补充。}

\item
%% number: 16
%% paperName: 2000年天津高考第 16 题 8 分
%% typeId: 4
%% type: 实验题
%% chapter: 电场
%% point: 示波器
%% method: 无
%% score: 8
%% degree: 550
%% duplicateId: 0
%% body:
图甲为示波器面板，图乙为一信号源。

\begin{enumerate}
	\item
	若要观测此信号源发出的正弦交流信号的波形，应将信号源的 a 端与示波器面板上的 \tkanswer{Y 输入} 接线柱相连，b 端与 \tkanswer{地} 接线柱相连。

	\item
	若示波器所显示的输入波形如图丙所示，要将波形上移，应调节面板上的 \tkanswer{竖直位移（填写 $\downarrow\uparrow$ 的也可）} 旋钮；要使此波形横向展宽，应调节 \tkanswer{X 增益} 旋钮；要使屏上能够显示 3 个完整的波形，应调节 \tkanswer{扫描范围和扫描微调} 旋钮。
\end{enumerate}
\onepicture[align=right, w=8cm]{figs/16.png}

%% answer: （1）Y 输入，地；（2）竖直位移，X 增益，扫描范围和扫描微调
\jdanswer{
\begin{enumerate}
	\item
	Y 输入，地

	\item
	竖直位移（填写 $\downarrow\uparrow$ 的也可），X 增益，扫描范围和扫描微调
\end{enumerate}
}

%% memo:
\memoanswer{本题原卷及材料中未提供详解，待补充。}

\item
%% number: 17
%% paperName: 2000年天津高考第 17 题 11 分
%% typeId: 6
%% type: 计算题
%% chapter: 气体性质
%% point: 克拉珀龙方程
%% method: 无
%% score: 11
%% degree: 650
%% duplicateId: 0
%% body:
有一实用氧气钢瓶，瓶内氧气的压强 $p = 50\times10^{3}\UPa$，温度 $t = 27\UdC$，求氧气的密度，氧的摩尔质量 $\mu = 3.2\times10^{2}\Ukgmol$，结果取两位数字。

%% answer: \rho = 64\Ukgmc
\jdanswer{$\rho = 64\Ukgmc$}

%% memo:
\memoanswer{
设钢瓶内氧气的摩尔数为 $n$，体积为 $V$，则有

$n = \frac{pV}{RT}$ ①

氧气密度

$\rho = \frac{n\mu}{V}$ ②

由①、②式联立得

$\rho = \frac{p\mu}{RT}$ ③

以题给数据代入得

$\rho = 64\Ukgmc$ ④

评分标准：本题 11 分，①、②两式各 1 分。③式 6 分（没有①、②两式而直接写出③式的给 8 分），结果④式 3 分。
}

\item
%% number: 19
%% paperName: 2000年天津高考第 19 题 13 分
%% typeId: 6
%% type: 计算题
%% chapter: 力物体的平衡
%% point: 力矩平衡
%% method: 无
%% score: 13
%% degree: 450
%% duplicateId: 0
%% body:
如图~\ref{2000天津19} 所示是用电动砂轮打磨工件的装置，砂轮的转轴通过图中 O 点垂直于纸面，AB 是一长度 $l = 0.60\Um$、质量 $m_{1} = 0.50\Ukg$ 的均匀刚性细杆，可绕过 A 端的固定轴在竖直面（图中纸面）内无摩擦地转动，工件 C 固定在 AB 杆上，其质量 $m_{2} = 1.5\Ukg$，工件的重心、工件与砂轮的接触点 P 以及 O 点都在过 AB 中点的竖直线上，P 到 AB 杆的垂直距离 $d = 0.1\Um$，AB 杆始终处于水平位置，砂轮与工件之间的动摩擦因数 $\mu = 0.6$。

\begin{enumerate}
	\item
	当砂轮静止时，要使工件对砂轮的压力 $F_{0} = 100\UN$，则施于 B 端竖直向下的力 $F_{B}$ 应是多大？

	\item
	当砂轮逆时针转动时，要使工件对砂轮的压力仍为 $F_{0} = 100\UN$，则施于 B 端竖直向下的力 $F_{B}'$ 应是多大？
\end{enumerate}
\onepicture[num=19, showlabel=false, label=2000天津19, align=right, w=5cm]{figs/19.png}

%% answer: （1）40\UN；（2）30\UN
\jdanswer{
\begin{enumerate}
	\item
	$40\UN$

	\item
	$30\UN$
\end{enumerate}
}

%% memo:
\memoanswer{本题原卷及材料中未提供详解，待补充。}

\end{enumerate}

\end{document}
